\capituloPregunta{¿Por qué necesitamos diseñar experimentos?}

\section{Motivación}

En ingeniería se toman decisiones con información incompleta. Se cambia un material, una temperatura, una herramienta o una mezcla y se observa un resultado. El problema aparece cuando no sabemos si el cambio observado se debe realmente al tratamiento o a otra condición que se modificó al mismo tiempo.

\begin{idea}
Diseñar un experimento significa construir una situación donde la evidencia pueda sostener una decisión. No se trata de probar cosas al azar, sino de controlar la forma en que se aprende de la realidad.
\end{idea}

\section{Intuición}

Supongamos que una mezcla produce un color mas intenso. Si también cambió el operador, la hora del día, el instrumento de medición y el lote de material, no se puede afirmar que la mezcla sea la causa del cambio. La observación puede ser útil, pero todavía no es evidencia experimental.

\section{Fundamento mínimo}

Un experimento requiere al menos:

\begin{itemize}
  \item una pregunta clara;
  \item una unidad experimental;
  \item una o más variables respuesta;
  \item factores que se modifican de forma controlada;
  \item condiciones que se mantienen constantes;
  \item repeticiones independientes;
  \item un criterio de interpretación.
\end{itemize}

\section{Ejemplo guiado}

Antes de calcular, el estudiante debe poder completar la siguiente tabla:

\begin{table}[H]
\centering
\begin{tabular}{p{0.28\textwidth}p{0.58\textwidth}}
\toprule
Elemento & Pregunta que responde\\
\midrule
Problema & ¿Qué decisión se quiere tomar?\\
Unidad experimental & ¿Sobre qué objeto se aplica el tratamiento?\\
Factor & ¿Qué condición se modifica?\\
Variable respuesta & ¿Qué se mide u observa?\\
Control & ¿Qué se mantiene constante?\\
Evidencia & ¿Qué dato permitirá decidir?\\
\bottomrule
\end{tabular}
\end{table}

\section{Errores frecuentes}

\begin{cuidado}
El error más común es confundir una experiencia llamativa con un experimento. Una observación puede iniciar una pregunta, pero no basta para concluir causalidad.
\end{cuidado}

\section{Conexión}

Una vez que la pregunta existe, el siguiente problema es aprender a medir y reconocer la variabilidad natural de los datos.
